% Trang 27 của Report.pdf. Chữ UTF-8 và nét vẽ TikZ có thể chỉnh sửa.
\begin{ReportPage}
\path[line width=0.398bp,line cap=butt,line join=miter,draw={rgb,1:red,0;green,0;blue,0},draw opacity=1] (70.866,73.90002) -- (530.07898,73.90002);
\path[line width=0.299bp,line cap=butt,line join=miter,draw={rgb,1:red,0;green,0;blue,0},draw opacity=1] (70.866,791.62) -- (530.07898,791.62);
\node[anchor=north west,inner sep=0pt] at (70.866,46.7751) {\includegraphics[width=26.66104bp,height=26.92589bp]{figures/original/image-263.png}};
% --- Khối văn bản ---
\node[anchor=base west,inner sep=0pt,outer sep=0pt] at (109.902,57.68604) {\ReportText{203.72517}{9.9626}{\ReportRoman}{0,0,0}{Trường Đại học Bách Khoa - ĐHQG TP.HCM}};
\node[anchor=base west,inner sep=0pt,outer sep=0pt] at (109.902,69.64203) {\ReportText{167.93946}{9.9626}{\ReportRoman}{0,0,0}{Khoa Khoa học và Kỹ thuật Máy tính}};
% --- Khối văn bản ---
\node[anchor=base west,inner sep=0pt,outer sep=0pt] at (99.213,98.83301) {\ReportText{430.84395}{10.9091}{\ReportRoman}{0,0,0}{Nhóm đề xuất cảnh báo khi trạm hoặc tuyến TCS mất phản hồi liên tục; lệnh UNKNOWN}};
\node[anchor=base west,inner sep=0pt,outer sep=0pt] at (70.866,114.41498) {\ReportText{459.14218}{10.9091}{\ReportRoman}{0,0,0}{còn tồn đọng sau 1 phút được đưa lên màn hình SEC. Truy vấn báo cáo được phân trang và giới}};
\node[anchor=base west,inner sep=0pt,outer sep=0pt] at (70.866,129.99597) {\ReportText{459.18588}{10.9091}{\ReportRoman}{0,0,0}{hạn tài nguyên, kèm chỉ mục theo nhân viên, thời gian và nhà máy để giữ kết nối DB cho xử lý}};
\node[anchor=base west,inner sep=0pt,outer sep=0pt] at (70.866,145.57794) {\ReportText{298.45111}{10.9091}{\ReportRoman}{0,0,0}{cổng. Cấu hình CPU/RAM sẽ được chọn sau các phép đo M2.}};
% --- Khối văn bản ---
\node[anchor=base west,inner sep=0pt,outer sep=0pt] at (70.866,175.79596) {\ReportText{297.97141}{11.9552}{\ReportBold}{0.11765,0.16078,0.23137}{8.6 Bước 6: kiểm tra giới hạn sẵn sàng và an toàn}};
% --- Khối văn bản ---
\node[anchor=base west,inner sep=0pt,outer sep=0pt] at (70.866,199.01196) {\ReportText{459.17502}{10.9091}{\ReportRoman}{0,0,0}{Với một ứng dụng và một DB, sự cố trung tâm có thể làm gián đoạn quét tự động. Phương án}};
\node[anchor=base west,inner sep=0pt,outer sep=0pt] at (70.866,214.59296) {\ReportText{459.14218}{10.9091}{\ReportRoman}{0,0,0}{ban đầu là khôi phục dịch vụ và để bảo vệ ghi nhận sự cố trong thời gian chờ. Nếu TMC yêu}};
\node[anchor=base west,inner sep=0pt,outer sep=0pt] at (70.866,230.17493) {\ReportText{459.15317}{10.9091}{\ReportRoman}{0,0,0}{cầu tiếp tục tự động khi mất WAN, nhóm cần bổ sung chương trình tại nhà máy có lưu cục bộ,}};
\node[anchor=base west,inner sep=0pt,outer sep=0pt] at (70.866,245.75592) {\ReportText{190.89835}{10.9091}{\ReportRoman}{0,0,0}{đồng bộ và kiểm tra quyền ngoại tuyến.}};
% --- Khối văn bản ---
\node[anchor=base west,inner sep=0pt,outer sep=0pt] at (99.213,264.60992) {\ReportText{430.78957}{10.9091}{\ReportRoman}{0,0,0}{Phân quyền theo nguyên tắc mặc định từ chối chỉ áp dụng cho truy cập thông thường. Khi}};
\node[anchor=base west,inner sep=0pt,outer sep=0pt] at (70.866,280.19092) {\ReportText{459.15311}{10.9091}{\ReportRoman}{0,0,0}{thoát hiểm, cháy, mất điện hoặc mở cơ khẩn cấp, cổng tuân theo cơ chế an toàn của nhà máy và}};
\node[anchor=base west,inner sep=0pt,outer sep=0pt] at (70.866,295.77289) {\ReportText{459.13138}{10.9091}{\ReportRoman}{0,0,0}{TCS. Dữ liệu những lượt này được đối soát sau để bảo đảm phần mềm chấm công không cản trở}};
\node[anchor=base west,inner sep=0pt,outer sep=0pt] at (70.866,311.35388) {\ReportText{54.29459}{10.9091}{\ReportRoman}{0,0,0}{việc sơ tán.}};
% --- Khối văn bản ---
\node[anchor=base west,inner sep=0pt,outer sep=0pt] at (70.866,350.51788) {\ReportText{196.4281}{14.3462}{\ReportBold}{0.11765,0.16078,0.23137}{9 Các quyết định kiến trúc}};
% --- Khối văn bản ---
\node[anchor=base west,inner sep=0pt,outer sep=0pt] at (70.866,378.23389) {\ReportText{459.13126}{10.9091}{\ReportRoman}{0,0,0}{Phần này ghi lại những lựa chọn chính, lý do chọn và ảnh hưởng của chúng tới hệ thống. Nhóm}};
\node[anchor=base west,inner sep=0pt,outer sep=0pt] at (70.866,393.81488) {\ReportText{459.10928}{10.9091}{\ReportRoman}{0,0,0}{dùng các quyết định này làm cơ sở cho M2. Các điểm phụ thuộc API, chính sách nhân sự và điều}};
\node[anchor=base west,inner sep=0pt,outer sep=0pt] at (70.866,409.39688) {\ReportText{246.12015}{10.9091}{\ReportRoman}{0,0,0}{kiện vận hành được tổng hợp trong sổ rủi ro ở Mục}};
\node[anchor=base west,inner sep=0pt,outer sep=0pt] at (316.98615,409.39688) {\ReportText{14.54184}{10.9091}{\ReportRoman}{0.03922,0.17255,0.65882}{ 11}};
\node[anchor=base west,inner sep=0pt,outer sep=0pt] at (331.52798,409.39688) {\ReportText{3.03271}{10.9091}{\ReportRoman}{0,0,0}{.}};
% --- Khối văn bản ---
\node[anchor=base west,inner sep=0pt,outer sep=0pt] at (70.866,439.61487) {\ReportText{330.78842}{11.9552}{\ReportBold}{0.11765,0.16078,0.23137}{9.1 ADR-01: Tổ chức ứng dụng theo mô-đun nghiệp vụ}};
% --- Khối văn bản ---
\node[anchor=base west,inner sep=0pt,outer sep=0pt] at (70.866,462.83087) {\ReportText{52.18914}{10.9091}{\ReportBold}{0,0,0}{Bối cảnh.}};
\node[anchor=base west,inner sep=0pt,outer sep=0pt] at (123.05514,462.83087) {\ReportText{406.98563}{10.9091}{\ReportRoman}{0,0,0}{ TMC có hai nhà máy với khoảng 1.400 nhân viên. Nhóm cần một cấu trúc dễ phân}};
\node[anchor=base west,inner sep=0pt,outer sep=0pt] at (70.866,478.41187) {\ReportText{270.32748}{10.9091}{\ReportRoman}{0,0,0}{công, kiểm thử và triển khai trong M2 (NFR-1, NFR-5).}};
% --- Khối văn bản ---
\node[anchor=base west,inner sep=0pt,outer sep=0pt] at (99.213,497.26587) {\ReportText{65.52005}{10.9091}{\ReportBold}{0,0,0}{Quyết định.}};
\node[anchor=base west,inner sep=0pt,outer sep=0pt] at (164.73305,497.26587) {\ReportText{365.29137}{10.9091}{\ReportRoman}{0,0,0}{ Sáu mô-đun được triển khai trong một ứng dụng. Mỗi mô-đun có giao diện}};
\node[anchor=base west,inner sep=0pt,outer sep=0pt] at (70.866,512.84686) {\ReportText{459.16385}{10.9091}{\ReportRoman}{0,0,0}{lập trình riêng, chia lớp bên trong và quản lý các bảng thuộc trách nhiệm của mình. Phần khởi}};
\node[anchor=base west,inner sep=0pt,outer sep=0pt] at (70.866,528.42883) {\ReportText{199.68011}{10.9091}{\ReportRoman}{0,0,0}{tạo đối tượng được tập trung tại một nơi.}};
% --- Khối văn bản ---
\node[anchor=base west,inner sep=0pt,outer sep=0pt] at (99.213,547.28186) {\ReportText{153.12018}{10.9091}{\ReportBold}{0,0,0}{Hệ quả và phương án khác.}};
\node[anchor=base west,inner sep=0pt,outer sep=0pt] at (252.33318,547.28186) {\ReportText{277.75648}{10.9091}{\ReportRoman}{0,0,0}{ Cách này giữ giao dịch và triển khai đơn giản, nhưng sự}};
\node[anchor=base west,inner sep=0pt,outer sep=0pt] at (70.866,562.86389) {\ReportText{459.10941}{10.9091}{\ReportRoman}{0,0,0}{cố tiến trình có thể ảnh hưởng toàn ứng dụng. Phương án phân lớp truyền thống vẫn phù hợp ở}};
\node[anchor=base west,inner sep=0pt,outer sep=0pt] at (70.866,578.44482) {\ReportText{459.1312}{10.9091}{\ReportRoman}{0,0,0}{mức cơ bản; tách thành microservices sẽ được xem xét khi có nhu cầu phát hành, mở rộng hoặc}};
\node[anchor=base west,inner sep=0pt,outer sep=0pt] at (70.866,594.02686) {\ReportText{119.98918}{10.9091}{\ReportRoman}{0,0,0}{vận hành độc lập rõ hơn.}};
% --- Khối văn bản ---
\node[anchor=base west,inner sep=0pt,outer sep=0pt] at (70.866,624.24487) {\ReportText{328.23007}{11.9552}{\ReportBold}{0.11765,0.16078,0.23137}{9.2 ADR-02: Dùng DB quan hệ và kiểm soát ghi trùng}};
% --- Khối văn bản ---
\node[anchor=base west,inner sep=0pt,outer sep=0pt] at (70.866,647.46088) {\ReportText{51.99278}{10.9091}{\ReportBold}{0,0,0}{Bối cảnh.}};
\node[anchor=base west,inner sep=0pt,outer sep=0pt] at (122.85878,647.46088) {\ReportText{407.17113}{10.9091}{\ReportRoman}{0,0,0}{ Yêu cầu từ kiosk hoặc thông báo từ TCS có thể được gửi nhiều lần. Dữ liệu công và}};
\node[anchor=base west,inner sep=0pt,outer sep=0pt] at (70.866,663.04187) {\ReportText{260.84747}{10.9091}{\ReportRoman}{0,0,0}{nhật ký cần được lưu nhất quán (FR-3.1–3.3, NFR-2).}};
% --- Khối văn bản ---
\node[anchor=base west,inner sep=0pt,outer sep=0pt] at (99.213,681.89587) {\ReportText{65.30186}{10.9091}{\ReportBold}{0,0,0}{Quyết định.}};
\node[anchor=base west,inner sep=0pt,outer sep=0pt] at (164.51486,681.89587) {\ReportText{365.50943}{10.9091}{\ReportRoman}{0,0,0}{ Nhóm chọn PostgreSQL, phân bảng theo mô-đun. Yêu cầu và nhật ký được}};
\node[anchor=base west,inner sep=0pt,outer sep=0pt] at (70.866,697.47687) {\ReportText{459.15323}{10.9091}{\ReportRoman}{0,0,0}{lưu cùng giao dịch; khi hoàn tất, kết quả và công cũng được lưu nguyên tử. Cơ chế chống trùng}};
\node[anchor=base west,inner sep=0pt,outer sep=0pt] at (70.866,713.05884) {\ReportText{459.22952}{10.9091}{\ReportRoman}{0,0,0}{kết hợp khóa UNIQUE, requestHash, CAS và trả lại trạng thái đã có. DB được cấu hình WAL,}};
\node[anchor=base west,inner sep=0pt,outer sep=0pt] at (70.866,728.63983) {\ReportText{176.34563}{10.9091}{\ReportRoman}{0,0,0}{ghi bền vững và sao lưu để phục hồi.}};
% --- Khối văn bản ---
\node[anchor=base west,inner sep=0pt,outer sep=0pt] at (99.213,747.49286) {\ReportText{154.40742}{10.9091}{\ReportBold}{0,0,0}{Hệ quả và phương án khác.}};
\node[anchor=base west,inner sep=0pt,outer sep=0pt] at (253.62042,747.49286) {\ReportText{276.44739}{10.9091}{\ReportRoman}{0,0,0}{ Một DB giúp các mô-đun chia sẻ giao dịch cục bộ. Tuy}};
\node[anchor=base west,inner sep=0pt,outer sep=0pt] at (70.866,763.07483) {\ReportText{459.15317}{10.9091}{\ReportRoman}{0,0,0}{nhiên, khóa duy nhất trong DB chỉ ngăn bản ghi trùng; việc ngăn phát lệnh mở lặp còn cần quy}};
% --- Khối văn bản ---
\node[anchor=base west,inner sep=0pt,outer sep=0pt] at (70.866,803.74701) {\ReportText{121.10019}{8.9664}{\ReportRoman}{0,0,0}{Kiến trúc phần mềm - HK261}};
\node[anchor=base west,inner sep=0pt,outer sep=0pt] at (493.97299,803.74701) {\ReportText{36.11661}{8.9664}{\ReportRoman}{0,0,0}{Trang 20}};
\end{ReportPage}
